\exercisesection{Homological Dimension}

\begin{enumerate}
\def\labelenumi{\arabic{enumi})}
\tightlist
\item
  Let \(u: \mathbf{A} \to \mathbf{B}\) be a ring homomorphism, M a B-module. Prove the inequality
\end{enumerate}

\[
\mathrm{dp} _ {\mathbf {A}} (\mathbf {M}) \leqslant \mathrm{dp} _ {\mathbf {B}} (\mathbf {M}) + \mathrm{dp} _ {\mathbf {A}} (\mathbf {B} _ {s}).
\]

(Reason by induction on dpB(M) or use the spectral sequence of Exercise 6, p.~188.)

\begin{enumerate}
\def\labelenumi{\arabic{enumi})}
\setcounter{enumi}{1}
\item
  Assume the ring A is noetherian; let M be a finitely generated A-module of projective dimension n. Show that the k-module Ext \(_{A}^{n}\) (M, A) is nonzero.
\item
  Let \(n\) be an integer \(\geqslant 1\) . Show that \(\mathrm{dh}(\mathbf{M}_n(\mathbf{A})) = \mathrm{dh}(\mathbf{A})\) (cf.~VIII, § 4).
\item
  Let M be an A-module. One calls the injective dimension of M, and denotes it by \(\mathrm{di}_{\mathbf{A}}(\mathbf{M})\) , the infimum in \(\overline{\mathbf{Z}}\) of the lengths of injective resolutions of M.
\end{enumerate}

If n is an integer \(\geqslant\) 0, show that the following conditions are equivalent:

\(\alpha)\mathrm{di}_{\mathbf{A}}(\mathbf{M})\leqslant n;\)

β) \(\operatorname{Ext}_{\mathbf{A}}^{r}(\mathbf{N},\mathbf{M}) = 0\) for every A-module N and every integer \(r > n\) ;

\(\gamma)\ \mathrm{Ext}_{\mathbf{A}}^{n+1}(\mathbf{N},\mathbf{M}) = 0\) for every A-module \(\mathbf{N}\) ;

δ) \(\operatorname{Ext}_{\mathbf{A}}^{n+1}(\mathbf{N},\mathbf{M})=0\) for every cyclic A-module N;

ε) for every exact sequence 0 → M → I⁰ → \(\cdots\) → Iⁿ⁻¹ → N → 0, where the Iᵏ are injective, N is injective. 5) Let (Mᵢ)ᵢ∈F be a family of A-modules.

\begin{enumerate}
\def\labelenumi{\alph{enumi})}
\item
  Show that \(\mathrm{di}_{\mathbf{A}}(\prod \mathbf{M}_i) = \sup \mathrm{di}_{\mathbf{A}}(\mathbf{M}_i)\) .
\item
  If A is noetherian, show that \(\mathrm{di}_{\mathbf{A}}(\bigoplus \mathbf{M}_i) = \sup \mathrm{di}_{\mathbf{A}}(\mathbf{M}_i)\) .
\item
  If \(\mathrm{di}_{\mathbf{A}}(\bigoplus_{i\in E}\mathbf{M}_{i}) = \sup_{i\notin E}\mathrm{di}_{\mathbf{A}}(\mathbf{M}_{i})\) for every family \((\mathbf{M}_i)_{i\in E}\) of A-modules, A is noetherian (use Exercise 21, p.~170).
\end{enumerate}

\begin{enumerate}
\def\labelenumi{\arabic{enumi})}
\setcounter{enumi}{5}
\item
  Let I be the set of left ideals of A. Show that \(\mathrm{dh}(\mathbf{A}) = \sup \mathrm{dp}_{\mathbf{A}}(\mathbf{A}/\mathfrak{a})\) .
\item
  Let M be an A-module. We call the flat dimension of M, and denote it by \(dpl_{A}\) (M), the infimum in \(\overline{Z}\) of the lengths of the flat resolutions of M.
\end{enumerate}

\begin{enumerate}
\def\labelenumi{\alph{enumi})}
\tightlist
\item
  If n is an integer \(\geqslant 0\) show that the following conditions are equivalent:
\end{enumerate}

\(\alpha)\mathrm{dpl}_{\mathbf{A}}(\mathbf{M})\leqslant n.\)

\(\beta)\ \mathrm{Tor}_{r}^{\mathbf{A}}(\mathbf{N},\mathbf{M})=0\ \text{for every right A-module N and every integer }r>n.\)

\(\gamma)\ \mathrm{Tor}_{n+1}^{\mathbf{A}}(\mathbf{N},\mathbf{M})=0\) for every right A-module N.

\(\delta)\ \mathrm{Tor}_{n+1}^{\mathbf{A}}(\mathbf{N},\mathbf{M})=0\) for every monogenic right A-module N.

ε) For every exact sequence \(0 \rightarrow K \rightarrow P_{n-1} \rightarrow \cdots \rightarrow P_{0} \rightarrow M \rightarrow 0\) where the \(P_{k}\) are flat, K is flat.

\begin{enumerate}
\def\labelenumi{\alph{enumi})}
\setcounter{enumi}{1}
\item
  Show that one has \(\mathrm{dpl}_{\mathbf{A}}(\mathbf{M}) \leqslant \mathrm{dp}_{\mathbf{A}}(\mathbf{M})\) ; there is equality if A is noetherian and M is finitely generated.
\item
  Let \((\mathbf{N}_{\alpha})_{\alpha \in I}\) be an inductive system of A-modules and N its inductive limit. Prove the inequality
\end{enumerate}

\[
\mathrm{dpl} _ {A} (N) \leqslant \sup _ {\alpha \in I} \mathrm{dpl} _ {A} (N _ {\alpha}).
\]

\begin{enumerate}
\def\labelenumi{\arabic{enumi})}
\setcounter{enumi}{7}
\tightlist
\item
  We call the tor-dimension of A and denote by td(A) the supremum in \(\overline{\mathbf{Z}}\) of the set of integers \(n\) for which there exist a left A-module N and a right A-module M such that \(\operatorname{Tor}_{n}^{\mathbf{A}}(\mathbf{M},\mathbf{N})\neq0\) .
\end{enumerate}

\begin{enumerate}
\def\labelenumi{\alph{enumi})}
\tightlist
\item
  Let \(n\) be an integer \(\geqslant 0\) . Show that the following conditions are equivalent:
\end{enumerate}

\(\alpha)\ \mathrm{td}(\mathbf{A}) \leqslant n.\)

β) For every A-module M, we have \(\mathrm{dpl}_{\mathbf{A}}(\mathbf{M}) \leqslant n\) (Exercise 7).

\(\gamma)\) For every monogenic A-module M, we have \(\mathrm{dpl}_{\mathbf{A}}(\mathbf{M}) \leqslant n\) .

\begin{enumerate}
\def\labelenumi{\alph{enumi})}
\setcounter{enumi}{1}
\tightlist
\item
  Show that we have td(A) = td(A°) and td(A) ≤ dh(A); equality holds if A is left noetherian.
\end{enumerate}

¶ 9) a) Show that the following conditions are equivalent:

α) Every finitely generated (left) ideal of A is a projective A-module.

β) Every finitely generated submodule of a projective A-module is projective.

\(\gamma)\) We have td(A) \(\leqslant\) 1 (Exercise 8), and the ring A is coherent (X, p.~181, Exercise 11).

δ) For every set I, every submodule of A¹ is flat.

\begin{enumerate}
\def\labelenumi{\alph{enumi})}
\setcounter{enumi}{1}
\tightlist
\item
  Suppose that the ring A satisfies the equivalent conditions of a). Prove that every projective A-module is isomorphic to a direct sum of finitely generated ideals of A. (First treat the case of a finitely generated module; then treat by induction the case of a module admitting a countable family of generators, observing that every element of a projective A-module P is contained in a finitely generated direct summand of P. Deduce the general case from Kaplansky's theorem (II, p.~183, Exercise 2).)
\end{enumerate}

\begin{enumerate}
\def\labelenumi{\arabic{enumi})}
\setcounter{enumi}{9}
\tightlist
\item
  Assume \(\mathbf{A}\) is an integral domain; let \(\mathbf{K}\) be its field of fractions. We say that an ideal \(\mathfrak{a}\) of \(\mathbf{A}\) is \(invertible\) if there exist elements \(x_{1},\ldots ,x_{n}\) of \(\mathbf{K}\) such that \(x_{i}\mathfrak{a}\subset \mathbf{A}\) for \(1\leqslant i\leqslant n\) and \(\sum x_i\mathfrak{a} = \mathbf{A}\) .
\end{enumerate}

\begin{enumerate}
\def\labelenumi{\alph{enumi})}
\item
  Show that every invertible ideal is finitely generated.
\item
  Show that a nonzero ideal of A is a projective A-module if and only if it is invertible.
\item
  Let a be an invertible ideal of A, D a divisible A-module. Prove that \(\operatorname{Ext}_A^i(A/a, D) = 0\) for \(i \geqslant 1\) .
\end{enumerate}

\begin{enumerate}
\def\labelenumi{\arabic{enumi})}
\setcounter{enumi}{10}
\tightlist
\item
  \begin{enumerate}
  \def\labelenumii{\alph{enumii})}
  \tightlist
  \item
    Suppose the ring A is an integral domain. Show that the conditions \(\alpha\) to \(\delta\) of Exercise 9 are still equivalent to the following conditions:
  \end{enumerate}
\end{enumerate}

ε) Every finitely generated torsion-free A-module is projective.

ζ) Every torsion-free A-module is flat.

η) Every submodule of a flat A-module is flat.

(Use Exercise 13, p.~169.)

An integral domain satisfying these conditions is said to be Prüferian.

\begin{enumerate}
\def\labelenumi{\alph{enumi})}
\setcounter{enumi}{1}
\tightlist
\item
  Let B be an integral domain, the union of an increasing filtered family of Prüferian subrings. Show that B is Prüferian.
\end{enumerate}

* c) Show that the ring of algebraic integers is Prüferian. *

\begin{enumerate}
\def\labelenumi{\arabic{enumi})}
\setcounter{enumi}{11}
\tightlist
\item
  Suppose the ring A is an integral domain. Show that the following conditions are equivalent:
\end{enumerate}

\(\alpha)\ \mathrm{dh}(\mathbf{A}) \leqslant 1.\)

β) Every divisible A-module is injective.

γ) The ring A is Prüferian (Exercise 11) and noetherian.

(To show that α) implies β), use Exercise 10.)

We then say that A is a Dedekind ring.

\begin{enumerate}
\def\labelenumi{\arabic{enumi})}
\setcounter{enumi}{12}
\tightlist
\item
  Assume A is left and right noetherian. Show that the following conditions are equivalent:
\end{enumerate}

\(\alpha)\ \mathrm{dh}(\mathbf{A}) \leqslant 2.\)

β) For all integers \(p, q\) and every A-linear map \(u: \mathbf{A}^p \to \mathbf{A}^q\) , the kernel of \(u\) is a projective A-module.

γ) The dual of every finitely generated A-module is projective.

\begin{enumerate}
\def\labelenumi{\arabic{enumi})}
\setcounter{enumi}{13}
\tightlist
\item
  Let x be an element of A that is left-cancellable.
\end{enumerate}

\begin{enumerate}
\def\labelenumi{\alph{enumi})}
\item
  Show that the ideal \(x\mathbf{A}\) is two-sided if and only if there exists an endomorphism \(\sigma\) of the ring \(\mathbf{A}\) such that \(ax = x\sigma(a)\) for all \(a \in \mathbf{A}\) .
\item
  We now assume that the conditions of a) are satisfied. Let M be an A-module such that xM = 0. Prove the equality dp\_A(M) = dp\_{A/xA}(M) + 1 (one may use the spectral sequence of Exercise 6, d), p. 188).
\item
  We assume that \(x\) belongs to the radical of \(\mathbf{A}\) , and that \(\mathbf{A}\) is noetherian. Let \(\mathbf{N}\) be a finitely generated \(\mathbf{A}\) -module, such that the homothety of ratio \(x\) is injective in \(\mathbf{N}\) . Show that we have \(\mathrm{dp}_{\mathbf{A}}(\mathbf{N}) = \mathrm{dp}_{\mathbf{A}/\mathbf{A}x}(\mathbf{N}/x\mathbf{N})\) . d) Under the conditions of \(c\) ), show that one has \(\mathrm{dh}(\mathbf{A}) = \mathrm{dh}(\mathbf{A}/\mathbf{A}x) + 1\) .
\end{enumerate}

\begin{enumerate}
\def\labelenumi{\arabic{enumi})}
\setcounter{enumi}{14}
\tightlist
\item
  Let σ be an automorphism of the ring A.
\end{enumerate}

\begin{enumerate}
\def\labelenumi{\alph{enumi})}
\tightlist
\item
  We denote by \(\mathbf{A}_{\sigma}[\mathbf{X}]\) the ring defined in VIII, § 1, no. 3. Show that
\end{enumerate}

\[
\mathrm{dh} (\mathbf {A} _ {\sigma} [ \mathbf {X} ]) = \mathrm{dh} (\mathbf {A}) + 1.
\]

(The inequality dh(A \(_{\sigma}\) {[}X{]}) \(\geqslant\) dh(A) + 1 follows from Exercise 14, \(b\) ); for the opposite inequality, be inspired by the proof of X, p.~142, lemma 3, \(a\) .)

\begin{enumerate}
\def\labelenumi{\alph{enumi})}
\setcounter{enumi}{1}
\tightlist
\item
  We denote by \(\mathbf{A}_{\sigma}[[\mathbf{X}]]\) the ring defined in IV, § 4, Exercise 8. Show that if A is noetherian, then we have \(\mathrm{dh}(\mathbf{A}_{\sigma}[[\mathbf{X}]))=\mathrm{dh}(\mathbf{A})+1\) (use Exercise 14).
\end{enumerate}

\begin{enumerate}
\def\labelenumi{\arabic{enumi})}
\setcounter{enumi}{15}
\tightlist
\item
  Let A be a local ring that is noetherian on the right and on the left, and let m be its maximal ideal; we denote by \(k_{s}\) (resp. \(k_{d}\) ) the left (resp. right) A-module A/m. Let M be a Tor \(_{n+1}^{A}\) finitely generated left module and let n be an integer. a) Show that one has dp \(_{A}\) (M) \(\leqslant\) n if and only if Tor \(_{n+1}^{A}\) ( \(k_{d}\) , M) = 0 (use Exercise 4, p.~185).
\end{enumerate}

\begin{enumerate}
\def\labelenumi{\alph{enumi})}
\setcounter{enumi}{1}
\tightlist
\item
  Show that one has \(\mathrm{dh(A)} = \mathrm{dp}_{\mathbf{A}}(k_s)\) ; for \(\mathrm{dh(A)} \leqslant n\) it is necessary and sufficient that one has
\end{enumerate}

\[
\operatorname{Tor} _ {n + 1} ^ {\mathbf {A}} \left(k _ {d}, k _ {s}\right) = 0.
\]

\begin{enumerate}
\def\labelenumi{\alph{enumi})}
\setcounter{enumi}{2}
\tightlist
\item
  Let \(x\) be an element of the center of \(\mathbf{A}\) belonging to \(m\) , such that the homothety with ratio \(x\) in \(M\) is injective. Show that we have \(\mathrm{dp}_{\mathrm{A}}(\mathrm{M}/x\mathrm{M}) = \mathrm{dp}_{\mathrm{A}}(\mathrm{M}) + 1\) (cf. also Exercise 14).
\end{enumerate}

\begin{enumerate}
\def\labelenumi{\arabic{enumi})}
\setcounter{enumi}{16}
\tightlist
\item
  Let M be an A-module, I a well-ordered set, \((\mathbf{M}_{\alpha})_{\alpha \in I}\) an increasing family of submodules of M, such that \(M = \bigcup M_{\alpha}\) . We set \(M_{\alpha}^{\prime} = \bigcup M_{\beta}\) for \(\alpha \in I\) .
\end{enumerate}

\begin{enumerate}
\def\labelenumi{\alph{enumi})}
\item
  Prove the inequality dpA(M) ≤ sup dpA(Mα/Mα').
\item
  Suppose there exists an integer \(n\) such that \(\mathrm{dp}_{\mathbf{A}}(\mathbf{M}_{\alpha}') \leqslant n\) for all \(\alpha \in I\) . Show that we have
\end{enumerate}

\[
\mathrm{dp} _ {\mathbf {A}} (\mathbf {M}) \leqslant n + 1.
\]

\begin{enumerate}
\def\labelenumi{\arabic{enumi})}
\setcounter{enumi}{17}
\tightlist
\item
  \begin{enumerate}
  \def\labelenumii{\alph{enumii})}
  \tightlist
  \item
    Let \((\mathbf{M}_{\alpha})_{\alpha \in I}\) be an inductive system of A-modules relative to a countable set I, M its limit, \(n\) be an integer \(\geqslant 0\) . Show that if \(\mathrm{dp}_{\mathbf{A}}(\mathbf{M}_{\alpha}) \leqslant n\) for all \(\alpha \in I\) , we have \(\mathrm{dp}_{\mathbf{A}}(\mathbf{M}) \leqslant n + 1\) (reduce to the case \(n = 0\) and \(I = N\) , then write an exact sequence \(0 \to \bigoplus_{n} \mathbf{M}_{n} \to \bigoplus_{n} \mathbf{M}_{n} \to \mathbf{M} \to 0\) ).
  \end{enumerate}
\end{enumerate}

\begin{enumerate}
\def\labelenumi{\alph{enumi})}
\setcounter{enumi}{1}
\item
  Let P be a flat A-module admitting a presentation \(\mathbf{A}^{(\mathbb{S})}\to\mathbf{A}^{(\mathbb{T})}\to\mathbf{P}\to0\) , where the set S is countable. Deduce from a) that one has dp \(_{\mathbb{A}}(\mathbb{P})\leqslant1\) .
\item
  Assume that every left ideal of the ring A is generated by a countable set of elements. Prove that dpA(M) ≤ dplA(M) + 1 for every A-module M generated by a countable family of elements. Deduce the inequalities dh(A) ≤ td(A) + 1 and \textbar dh(A) - dh(A°)\textbar{} ≤ 1. (Use Exercise 12, p.~181, Exercise 8, p.~202, as well as b).)
\end{enumerate}

¶ 19) a) Let n be an integer ≥ 0, let (Mα)α ∈ I be an inductive system of A-modules, and let M be its inductive limit. Assume that Card (I) ≤ Nn (E, III, p.~87, exercise 10) and that there exists an integer r such that dpA(Mα) ≤ r for every α ∈ I. Show that dpA(M) ≤ r + n + 1 (argue by induction on n, using exercise 18, a) and exercise 17).

\begin{enumerate}
\def\labelenumi{\alph{enumi})}
\setcounter{enumi}{1}
\tightlist
\item
  We assume that every left ideal of A is generated by a family of elements of cardinal \(\leqslant \aleph_{n}\) . Prove that one has dp \(_{A}\) (M) \(\leqslant\) dpl \(_{A}\) (M) + n + 1 for every A-module M generated by a family of elements of cardinality \(\leqslant \aleph_{n}\) . Deduce from this the inequalities
\end{enumerate}

\[
\mathrm{dh} (\mathbf {A}) \leqslant \mathrm{td} (\mathbf {A}) + n + 1 \quad \text { and } \quad | \mathrm{dh} (\mathbf {A}) - \mathrm{dh} (\dot {\mathbf {A}} ^ {\circ}) | \leqslant n + 1.
\]

¶ 20) Let R be a principal ideal ring, K its field of fractions. We denote by B the subring of M₂(K) consisting of the matrices \(\begin{pmatrix} a & 0 \\ x & y \end{pmatrix}\) with \(a \in \mathbb{R}, x, y \in \mathbb{K}\) .

\begin{enumerate}
\def\labelenumi{\alph{enumi})}
\item
  Show that B is left noetherian but not right noetherian.
\item
  Show that \(\mathrm{dh(B)} = 1\) and \(\mathrm{dh(B^{\circ})} = 2\) .
\end{enumerate}

\begin{enumerate}
\def\labelenumi{\arabic{enumi})}
\setcounter{enumi}{20}
\tightlist
\item
  Let r be the radical of A. Show that the following conditions are equivalent:
\end{enumerate}

α) A/r is semisimple, and for every sequence \((a_{n})_{n\geq1}\) of elements of r, there exists an integer n such that \(a_{1}\ldots a_{n}=0\) .

β) Every A-module admits a projective cover.

\(\gamma)\) For every A-module M, we have \(\mathrm{dpl}_{\mathbf{A}}(\mathbf{M}) = \mathrm{dp}_{\mathbf{A}}(\mathbf{M})(\mathbf{X}, p. 202, \text{exercise } 7)\) .

δ) For every inductive system \((\mathbf{M}_{\alpha}, f_{\beta\alpha})_{\alpha \in I}\) of A-modules, we have:

\[
\mathrm{dp} _ {\mathrm{A}} (\varinjlim \mathrm{M} _ {\alpha}) \leqslant \sup _ {\alpha \in I} \mathrm{dp} _ {\mathrm{A}} (\mathrm{M} _ {\alpha}).
\]

ε) Every flat A-module is projective.

φ) Every descending sequence of monogenic right ideals of A is stationary.

(For the equivalence of \(\alpha\) ) and \(\beta\) ), cf.~VIII, § 8, exercise 26. To show that \(\beta\) ) implies \(\gamma\) ), note that M admits a minimal projective resolution P and that \(P_n/rP_n\) is isomorphic to Tor \(_n^A\) (A/r, M). To prove \(\varepsilon) \Rightarrow \varphi\) ), let \((l_n)_{n \geq 1}\) be a decreasing sequence of monogenic right ideals, such that

\[
\mathrm{I} _ {n} = a _ {1} \dots a _ {n} \mathrm{A}, \quad \text { with } \quad a _ {i} \in \mathrm{A}.
\]

In the free module \(\mathbf{A}^{(\mathbf{N})}\) , consider the submodules \(\mathbf{N}_p\) generated by \(e_1 - a_1 e_2, \ldots, e_p - a_p e_{p+1}\) . Deduce from \(\varepsilon\) ) that \(\mathbf{N} = \bigcup \mathbf{N}_p\) is a direct summand in \(\mathbf{A}^{(\mathbf{N})}\) ; conclude from this that \(\mathrm{l}_{n+1} = \mathrm{l}_n\) for \(n\) large enough.

For the implication \(\varphi) \Rightarrow \alpha)\) , use VIII, § 9, exercise 12.)

\begin{enumerate}
\def\labelenumi{\arabic{enumi})}
\setcounter{enumi}{21}
\tightlist
\item
  Show that the following conditions are equivalent:
\end{enumerate}

(α) Every finitely presented A-module of finite projective dimension is projective.

(β) Every injective A-homomorphism \(A^{p} \rightarrow A^{q}\) admits a retraction.

(γ) For every finitely generated right ideal a distinct from A, there exists a nonzero element x of A such that xa = 0.

(δ) Every simple right A-module of finite presentation is isomorphic to a (minimal) right ideal of A.

(Observe that condition \(\gamma\) ) is equivalent to condition \(\beta\) in the case \(p = 1\) .)

¶ 23) Show that the following conditions are equivalent:

(α) Every A-module of finite projective dimension is projective.

(β) The ring A satisfies the equivalent conditions of Exercise 21, and for every finitely generated right ideal a distinct from A, there exists an element \(x \neq 0\) of A such that \(xa = 0\) .

(γ) The ring A satisfies the equivalent conditions of Exercise 21, and every simple A-module is a quotient of an injective A-module.

(Show that \(\alpha\) ) implies that every descending sequence of monogenic right ideals of A is stationary, by adapting the proof of Exercise 21. Then deduce from Exercise 22 the equivalence of \(\alpha\) ) and \(\beta\) ). To see that \(\alpha\) ) implies \(\gamma\) ), consider a projective cover P of the simple module M, an injective module I containing P, and a projective cover Q of I; define an A-homomorphism from Q to P, and deduce from it a surjection from I to M. To show that \(\gamma\) ) implies \(\alpha\) ), show that an A-module M such that dp \(_{A}\) (M) \(\leqslant\) 1 is projective by considering a projective cover \(p: P \to M\) and a surjection from Ker \(p\) onto a simple module.)

\begin{enumerate}
\def\labelenumi{\arabic{enumi})}
\setcounter{enumi}{23}
\tightlist
\item
  We denote by dhf(A) (resp. dif(A), resp. tdf(A)) the supremum of the projective (resp. injective, resp. flat) dimensions of A-modules for which this dimension is finite. If A has finite homological dimension, then dhf(A) = dif(A) = dh(A) and tdf(A) = td(A). If td(A) \textless{} ∞, we have tdf(A) = td(A). a) Show that tdf(A) ≤ dif(A°); equality holds if A is left noetherian (use X, p.~189, exercise 3).
\end{enumerate}

\begin{enumerate}
\def\labelenumi{\alph{enumi})}
\setcounter{enumi}{1}
\item
  Assume A is left noetherian. Show that we have dhf(A) ≤ di\_A(A\_s). If di\_A(A\_s) and dif(A) are finite, show that they are equal (use Exercise 2, p.~202).
\item
  Let G be a finite group not reduced to the identity element, and let A be the ring \(\mathbf{Z}^{(G)}\) . Show that we have
\end{enumerate}

\[
\mathrm{dh} (\mathbf {A}) = + \infty , \mathrm{dhf} (\mathbf {A}) = \mathrm{dif} (\mathbf {A}) = \mathrm{tdf} (\mathbf {A}) = 1, \mathrm{di} _ {\mathbf {A}} (\mathbf {A} _ {s}) = + \infty
\]

(cf. exercise 20, p. 193).

\begin{enumerate}
\def\labelenumi{\alph{enumi})}
\setcounter{enumi}{3}
\tightlist
\item
  Assume the ring A is self-injective (cf. X, p. 171, exercise 26), not semisimple. Show that every non-projective A-module has infinite projective dimension and infinite injective dimension. Deduce that one has dh(A) = +∞, dhf(A) = dif(A) = tdf(A) = di\_A(A\_s) = 0.
\end{enumerate}

\begin{enumerate}
\def\labelenumi{\arabic{enumi})}
\setcounter{enumi}{24}
\tightlist
\item
  We denote by \(\mathbf{A}^e\) the \(k\) -algebra \(\mathbf{A} \otimes_k \mathbf{A}^*\) and we consider \(\mathbf{A}\) as a left \(\mathbf{A}^e\) -module. We set
\end{enumerate}

\[
\mathrm{da} _ {k} (\mathrm{A}) = \mathrm{dp} _ {\mathrm{A} ^ {e}} (\mathrm{A}).
\]

If n is an integer \(\geqslant0\) , the inequality \(\mathrm{da}_{k}(\mathrm{A})\leqslant n\) is therefore equivalent to the vanishing of \(\mathrm{H}^{n+1}(\mathrm{A},\mathrm{M})\) for every (A, A)-bimodule M (X, p.~195, exercise 26). One has \(\mathrm{da}_{k}(\mathrm{A})=\mathrm{da}_{k}(\mathrm{A}^{\circ})\) .

\begin{enumerate}
\def\labelenumi{\alph{enumi})}
\item
  Show that one has \(\mathrm{da}_k(\mathbf{M}_n(k)) = 0\) and \(\mathrm{da}_k(k[\mathbf{X}_1, \ldots, \mathbf{X}_n]) = n\) for all \(n \geqslant 0\) .
\item
  If \(k\) is a field, one has \(\mathrm{da}_{k}(\mathbf{A}) = 0\) if and only if the \(k\) -algebra \(\mathbf{A}\) is absolutely semisimple (cf.~VIII, § 11, n° 3, th. 2).
\item
  Let M be a free \(k\) -module. Prove that one has \(\mathrm{da}_{k}(\mathbf{T}_{k}(\mathbf{M})) = 1\) (If \((e_{i})_{i \in I}\) is a basis of M, prove that the elements \((e_{i} \otimes 1 - 1 \otimes e_{i})\) form a basis of the left ideal of \(\mathbf{A}^{e}\) kernel of multiplication \(\mathbf{A} \otimes_{k} \mathbf{A}^{-} \to \mathbf{A}\) , with \(\mathbf{A} = \mathbf{T}_{k}(\mathbf{M})\) ).
\item
  We now assume that the \(k\) -module A is projective. Let \(\rho: k \to k'\) be a homomorphism of commutative rings; show that one has \(\mathrm{da}_k(\mathbf{A} \otimes_k k') \leqslant \mathrm{da}_k(\mathbf{A})\) . If \(\rho\) is injective and if \(\rho(k)\) is a direct summand of the \(k\) -module \(k'\) , one has equality (use X, p.~195, exercise 26, c)).
\item
  Let B be a second k-algebra (associative and unital), projective over k. Show that
\end{enumerate}

\[
\mathrm{da} _ {k} (\mathbf {A} \otimes_ {k} \mathbf {B}) \leqslant \mathrm{da} _ {k} (\mathbf {A}) + \mathrm{da} _ {k} (\mathbf {B}).
\]

If moreover k is a field and A and B are finite-dimensional over k, then equality holds (use X, p.~195, exercise 26, d)).

\begin{enumerate}
\def\labelenumi{\alph{enumi})}
\setcounter{enumi}{5}
\tightlist
\item
  Assume that k is a field. Prove the inequalities \(\mathrm{dh}(\mathbf{A}) \leqslant \mathrm{da}_{k}(\mathbf{A})\) and \(\mathrm{dh}(\mathbf{A}^{\circ}) \leqslant \mathrm{da}_{k}(\mathbf{A})\) (cf. X, p. 196, exercise 26, e)).
\end{enumerate}

\begin{enumerate}
\def\labelenumi{\arabic{enumi})}
\setcounter{enumi}{25}
\tightlist
\item
  Let A be an augmented \(k\) -algebra (X, p.~196, exercise 27); one assumes that there exists a homomorphism \(\rho: \mathrm{A} \to \mathrm{A}^e\) which satisfies the conditions of X, p.~196, exercise 27, c). Show that \(\mathrm{da}_k(\mathbf{A}) = \mathrm{dp}_\mathbf{A}(k)\) ; if moreover \(k\) is a field, one has \(\mathrm{da}_k(\mathbf{A}) = \mathrm{dp}_\mathbf{A}(k) = \mathrm{dh}(\mathbf{A}) = \mathrm{dh}(\mathbf{A}^\circ)\) . (Use Exercise 27, b) and c), p.~196, and Exercise 26, e), p.~196.) In particular, if V is a nonzero \(k\) -vector space, we have \(\mathrm{dh}\,\mathbf{T}_k(\mathbf{V}) = 1\) .
\end{enumerate}

* 27) We assume that \(k\) is a field. Let \(g\) be a Lie algebra over \(k\) of dimension \(n\) , U its enveloping algebra. Show that \(dhU = n\) . (Use Exercise 26, Exercise 27, d), p.~196, and Exercise 28, c), p.~197.) *

¶ 28) Let G be a group. One calls the cohomological dimension of G, denoted dc(G), the projective dimension of the \(\mathbf{Z}^{(G)}\) -module Z. If \(n\) is an integer, we thus have dc(G) \(\leqslant n\) if and only if H \(^{q}\) (G, M)=0 for every \(\mathbf{Z}^{(G)}\) -module M and every \(q > n\) .

\begin{enumerate}
\def\labelenumi{\alph{enumi})}
\item
  Prove that one has \(\mathrm{da}_{\mathbf{Z}}(\mathbf{Z}^{(G)}) = \mathrm{dc}(G)\) (use Exercise 26 and Exercise 27, d), p.~196).
\item
  If G is a free group not reduced to the identity element, then dc(G) = 1.
\item
  Let H be a subgroup of G; prove the inequality dc(H) ≤ dc(G). If dc(G) \textless{} ∞ and if H has finite index in G, show that dc(H) = dc(G) (show that for q = dc(G), the homomorphism
\end{enumerate}

\[
t ^ {q}: \mathrm{H} ^ {q} (\mathrm{H}, \mathrm{M}) \rightarrow \mathrm{H} ^ {q} (\mathrm{G}, \mathrm{M})
\]

of Exercise 14, p.~191, is surjective).

\begin{enumerate}
\def\labelenumi{\alph{enumi})}
\setcounter{enumi}{3}
\item
  Show that the cohomological dimension of a finite group not reduced to the identity element is infinite. Deduce that if dc(G) \textless{} ∞, every element of G other than the identity element has infinite order (one then says that G is torsion-free).
\item
  If G is torsion-free and if H is a subgroup of finite index in G, show that dc(H) = dc(G). (If (P, p) is a projective resolution of \(\mathbf{Z}^{(H)}\) -module Z and if (G : H) = n, define on the graded \(\mathbf{Z}^{(H)}\) -module \(\mathbf{T}_{\mathbf{Z}}^{n}(\mathbf{P})\) a structure of \(\mathbf{Z}^{(G)}\) -complex so that \((\mathbf{T}_{\mathbf{Z}}^{n}(\mathbf{P}), p^{\otimes n})\) is a projective resolution of the \(\mathbf{Z}^{(G)}\) -module Z.)
\item
  If H is a normal subgroup of G, show that we have
\end{enumerate}

\[
\mathrm{dc} (\mathbf {G}) \leqslant \mathrm{dc} (\mathbf {H}) + \mathrm{dc} (\mathbf {G} / \mathbf {H}).
\]

